\chapter*{Abstract}
\addcontentsline{toc}{chapter}{Abstract}

Large Language Models (LLMs) are being increasingly deployed to safety-critical
and enterprise domains to generated structured data, yet they natively struggle
to adhere to complex, context-sensitive operational semantics. While existing
constrained decoding techniques successfully enforce context-free grammars using
Finite State Machines, they fail to guarantee richly typed semantic invariants
without incurring some large or otherwise prohibitive inference latencies. To bridge this
gap between probabilistic generation and strict formal verification, this paper
proposes a dual-objective neuro-symbolic pipeline and architecture.

First, we formalize a "Syntax Oracle" as a deterministic, abstract machine.
The Syntax Oracle acts as a strict runtime
failsafe, continuously maintaining the expected generation state. It applies
dynamic logit masking only when the neural model deviates from valid
specification paths, thereby guaranteeing at least 100\% formal structural
complicance.

Second, to maximize the model's efficacy and minimize use of the Syntax Oracle,
we introduce a semantic internalization pipeline. By utilizing the Syntax Oracle
offline to procedurally generate synthetic data enriched with reasoning traces,
we can apply Low-Rank Adaptation (LoRA) via Supervised Fine-Tuning. This would,
in turn, train the LLM to natively internalize the underlying context-sensitive
logic, rather than merely memorizing structural distributions.

Finally, the efficacy of the system is evaluated using a novel,
entropy-normalized "Complexity-Adjusted Mask Intervention Rate", alongside
standard latency and throughput benchmarks. Ultimately, this dual-approach aims
to demonstrate that parameter-efficient fine-tuning on structural reasoning can
drastically reduce a model's reliance on deterministic fallbacks, providing a
scalable, low-latency framework for deploying fully verified and correct LLMs.
